%==============================================================
%  lecons/analyse/tchebychev.tex — Polynômes de Tchebychev
%==============================================================

\lecontitle{Polynômes de Tchebychev}{Analyse — Approximation et polynômes trigonométriques}

%=============================================================
%  § 1 — DÉFINITION ET PREMIÈRES VALEURS
%=============================================================
\secnum{navyblue}{1}{Définition et premières valeurs}

\begin{tcolorbox}[thmbox]
  \textbf{Définition.}\enspace\index{polynôme de Tchebychev}
  Le \emph{$n$-ième polynôme de Tchebychev de première espèce} est l'unique
  polynôme $T_n \in \mathbb{R}[X]$ vérifiant
  \[
    T_n(\cos\theta) = \cos(n\theta) \qquad \forall\,\theta \in \mathbb{R}.
  \]

  \medskip
  \textbf{Récurrence.}\enspace\index{récurrence de Tchebychev}
  Les polynômes $T_n$ sont bien définis et vérifient
  \[
    T_0 = 1,\qquad T_1 = X,\qquad
    T_{n+1} = 2X\,T_n - T_{n-1} \quad (n \geq 1).
  \]
\end{tcolorbox}

\rubriq{Existence et unicité}

\begin{mdframed}[style=proofbar]
  \textit{Existence.} Par récurrence sur $n$ à partir de la formule
  d'addition de cosinus :
  \[
    \cos((n+1)\theta) = 2\cos\theta\cdot\cos(n\theta) - \cos((n-1)\theta).
  \]
  Posant $x = \cos\theta$, si $T_n(x) = \cos(n\theta)$ et
  $T_{n-1}(x) = \cos((n-1)\theta)$, alors
  $T_{n+1}(x) = 2x\,T_n(x) - T_{n-1}(x)$ est bien un polynôme en $x$.

  \smallskip
  \textit{Unicité.} Deux polynômes coïncidant sur $[-1,1]$ (qui contient
  une infinité de points) sont égaux. \hfill$\blacksquare$
\end{mdframed}

\rubriq{Premières valeurs}

\begin{mdframed}[style=exbar]
  \begin{align*}
    T_0(x) &= 1 \\
    T_1(x) &= x \\
    T_2(x) &= 2x^2 - 1 \\
    T_3(x) &= 4x^3 - 3x \\
    T_4(x) &= 8x^4 - 8x^2 + 1 \\
    T_5(x) &= 16x^5 - 20x^3 + 5x
  \end{align*}
  Le coefficient dominant de $T_n$ est $2^{n-1}$ pour $n \geq 1$.
\end{mdframed}

\decsep{}

%=============================================================
%  § 2 — PROPRIÉTÉS ALGÉBRIQUES
%=============================================================
\secnum{navyblue}{2}{Propriétés algébriques}

\begin{tcolorbox}[thmbox]
  \textbf{Proposition.}\enspace
  Pour tout $n \geq 0$ :
  \begin{enumerate}
    \item \textbf{Parité.}\enspace $T_n(-x) = (-1)^n T_n(x)$.
    \item \textbf{Racines.}\enspace $T_n$ a exactement $n$ racines simples
          dans $(-1,1)$ :
          \[
            x_k = \cos\!\left(\frac{(2k-1)\pi}{2n}\right),
            \quad k = 1, \ldots, n.
          \]
    \item \textbf{Extrema.}\enspace $T_n$ prend alternativement les valeurs
          $+1$ et $-1$ aux $n+1$ points
          \[
            \xi_j = \cos\!\left(\frac{j\pi}{n}\right),
            \quad j = 0, 1, \ldots, n.
          \]
    \item \textbf{Borne.}\enspace $|T_n(x)| \leq 1$ pour tout $x \in [-1,1]$.
    \item \textbf{Composition.}\enspace $T_m \circ T_n = T_{mn}$.
  \end{enumerate}
\end{tcolorbox}

\rubriq{Preuve des points 1–3}

\begin{mdframed}[style=proofbar]
  \textit{Parité.} $T_n(-\cos\theta)=T_n(\cos(\pi-\theta))=\cos(n(\pi-\theta))
  =(-1)^n\cos(n\theta)=(-1)^n T_n(\cos\theta)$. \hfill$\square$

  \smallskip
  \textit{Racines.} $T_n(x_k)=\cos\!\bigl(n\cdot\tfrac{(2k-1)\pi}{2n}\bigr)
  =\cos\!\bigl(\tfrac{(2k-1)\pi}{2}\bigr)=0$.
  Ce sont $n$ valeurs distinctes dans $(-1,1)$, donc toutes les racines
  d'un polynôme de degré $n$. \hfill$\square$

  \smallskip
  \textit{Extrema.} $T_n(\xi_j)=\cos(j\pi)=(-1)^j$. \hfill$\blacksquare$
\end{mdframed}

\decsep{}

%=============================================================
%  § 3 — THÉORÈME D'APPROXIMATION MINIMALE
%=============================================================
\secnum{navyblue}{3}{Théorème d'approximation minimale}

\begin{tcolorbox}[thmbox]
  \textbf{Théorème (Tchebychev).}\enspace%
  \index{théorème d'approximation minimale de Tchebychev}
  Parmi tous les polynômes unitaires\footnote{%
    Coefficient dominant égal à $1$.}
  de degré $n \geq 1$, le polynôme
  \[
    \widetilde{T}_n = \frac{1}{2^{n-1}}\,T_n
  \]
  est celui qui minimise la norme uniforme sur $[-1,1]$ :
  \[
    \min_{\substack{P \in \mathbb{R}[X] \\ \deg P = n,\; P \text{ unitaire}}}
    \|P\|_{[-1,1]}
    \;=\;
    \|\widetilde{T}_n\|_{[-1,1]}
    \;=\;
    \frac{1}{2^{n-1}}.
  \]
\end{tcolorbox}

\rubriq{Preuve}

\begin{mdframed}[style=proofbar]
  On sait que $\|\widetilde{T}_n\|_{[-1,1]} = \frac{1}{2^{n-1}}$ (les $n+1$
  extrema alternent entre $\pm\frac{1}{2^{n-1}}$).

  \smallskip
  Soit $P$ un polynôme unitaire de degré $n$ tel que $\|P\|_{[-1,1]} <
  \frac{1}{2^{n-1}}$.
  Posons $Q = \widetilde{T}_n - P$ ; c'est un polynôme de degré $< n$
  (les termes dominants s'annulent). Aux $n+1$ points $\xi_j$~:
  \[
    \widetilde{T}_n(\xi_j) = \frac{(-1)^j}{2^{n-1}}, \qquad
    |P(\xi_j)| < \frac{1}{2^{n-1}},
  \]
  donc $Q(\xi_j)$ change de signe entre $\xi_{j-1}$ et $\xi_j$ ($n$ fois).
  Par le théorème des valeurs intermédiaires, $Q$ admet au moins $n$
  racines. Contradiction avec $\deg Q < n$, sauf si $Q = 0$, mais alors
  $P = \widetilde{T}_n$ et l'hypothèse est fausse. \hfill$\blacksquare$
\end{mdframed}

\decsep{}

%=============================================================
%  § 4 — ORTHOGONALITÉ
%=============================================================
\secnum{navyblue}{4}{Orthogonalité}

\begin{tcolorbox}[thmbox]
  \textbf{Théorème.}\enspace\index{orthogonalité des polynômes de Tchebychev}
  Les polynômes de Tchebychev sont orthogonaux pour le produit scalaire à
  poids $w(x)=(1-x^2)^{-1/2}$ sur $(-1,1)$~:
  \[
    \int_{-1}^{1} T_m(x)\,T_n(x)\,\frac{\mathrm{d}x}{\sqrt{1-x^2}}
    =
    \begin{cases}
      0       & \text{si } m \neq n, \\
      \pi     & \text{si } m = n = 0, \\
      \pi/2   & \text{si } m = n \geq 1.
    \end{cases}
  \]
\end{tcolorbox}

\rubriq{Preuve}

\begin{mdframed}[style=proofbar]
  On pose $x = \cos\theta$, d'où $\mathrm{d}x = -\sin\theta\,\mathrm{d}\theta$
  et $(1-x^2)^{1/2} = \sin\theta$.
  \[
    \int_{-1}^{1} T_m T_n \frac{\mathrm{d}x}{\sqrt{1-x^2}}
    = \int_0^{\pi} \cos(m\theta)\cos(n\theta)\,\mathrm{d}\theta.
  \]
  La formule de linéarisation $2\cos(m\theta)\cos(n\theta)
  = \cos((m-n)\theta)+\cos((m+n)\theta)$ et le calcul standard des intégrales
  de cosinus donnent le résultat. \hfill$\blacksquare$
\end{mdframed}

\decsep{}

%=============================================================
%  § 5 — APPLICATIONS, PIÈGES, EXERCICES
%=============================================================
\secnum{exgreen}{5}{Applications, pièges et exercices}

\noindent{\bfseries\color{navyblue} Applications.}

\begin{mdframed}[style=exbar]
  \textbf{1. Nœuds de Tchebychev pour l'interpolation.}\enspace%
  \index{nœuds de Tchebychev}
  Pour interpoler une fonction $f$ en $n+1$ points de $[-1,1]$, les racines
  de $T_{n+1}$ (nœuds de Tchebychev) minimisent le terme d'erreur, qui
  contient $\|\omega\|_\infty$ avec $\omega=\prod(x-x_k)$.

  \smallskip\noindent
  \textbf{2. Polynômes de Tchebychev de deuxième espèce.}\enspace%
  \index{polynôme de Tchebychev!de deuxième espèce}
  On définit $U_n$ par $U_n(\cos\theta)\sin\theta = \sin((n+1)\theta)$.
  Ils vérifient la même récurrence $U_{n+1}=2X\,U_n - U_{n-1}$ avec
  $U_0=1$, $U_1=2X$, et $T_n' = n\,U_{n-1}$.

  \smallskip\noindent
  \textbf{3. Développement en série de Tchebychev.}\enspace
  Toute fonction continue $f:[-1,1]\to\mathbb{R}$ admet un développement
  $f = \sum_{n\geq 0} c_n T_n$ convergent, analogue de la série de Fourier,
  avec $c_n = \frac{2}{\pi}\int_{-1}^1 f(x)T_n(x)\frac{\mathrm{d}x}{\sqrt{1-x^2}}$.
\end{mdframed}

\vspace{6pt}
\noindent{\bfseries\color{counterred} Pièges.}

\begin{mdframed}[style=cexbar]
  \textbf{Coefficient dominant $2^{n-1}$, pas $1$.}\enspace
  $T_n$ n'est pas unitaire pour $n\geq 2$. Le polynôme unitaire associé est
  $\widetilde{T}_n = \frac{1}{2^{n-1}}T_n$.

  \smallskip\noindent
  \textbf{Poids de l'orthogonalité.}\enspace
  L'orthogonalité de $(T_n)$ est relative au poids
  $(1-x^2)^{-1/2}$, non à la mesure de Lebesgue. Sans ce poids,
  $\int_{-1}^1 T_m T_n\,\mathrm{d}x \neq 0$ en général.

  \smallskip\noindent
  \textbf{Ne pas confondre $T_n$ et $U_n$.}\enspace
  $T_n' = n\,U_{n-1}$ ; les polynômes de deuxième espèce $U_n$ ont les
  racines $\cos\!\bigl(\frac{k\pi}{n+1}\bigr)$, $k=1,\ldots,n$, et non
  les racines de $T_n$.
\end{mdframed}

\vspace{6pt}
\exolabel

\begin{mdframed}[style=exobar]
  \begin{enumerate}
    \item Vérifier par la récurrence que $T_3(x) = 4x^3 - 3x$ et
          $T_4(x) = 8x^4 - 8x^2 + 1$, puis retrouver les racines de $T_3$
          via la formule trigonométrique.

    \item Montrer que $T_m(T_n(x)) = T_{mn}(x)$ en utilisant la définition
          $T_n(\cos\theta)=\cos(n\theta)$.

    \item En utilisant le théorème d'approximation minimale, montrer que
          pour tout polynôme $P$ de degré $\leq n-1$ :
          \[
            \max_{x\in[-1,1]}\left|\,x^n - P(x)\right| \geq \frac{1}{2^{n-1}}.
          \]

    \item Calculer $\int_{-1}^1 T_3(x)^2\,\frac{\mathrm{d}x}{\sqrt{1-x^2}}$
          directement, puis via le théorème d'orthogonalité.

    \item Montrer que $T_n'(x) = n\,U_{n-1}(x)$ en dérivant la relation
          $\cos(n\theta) = T_n(\cos\theta)$ par rapport à $\theta$.
  \end{enumerate}
\end{mdframed}

\leconfooter{P.\ L.\ Tchebychev, \textit{Théorie des mécanismes} (1853)}
